% 2012 年北京卷文科第 7 题：某三棱锥的三视图（正视图、侧视图、俯视图）。
% 来源 PDF：普通高考/2012/2012北京文.pdf，第 1 页（印面 719）中栏。
% 正视图：底 5（左段 2、右段 3），高 4；侧视图：底 4、高 4 的直角三角形；
% 俯视图：直角边 5 与 4 的直角三角形，顶点投影落在 5 的边上距左端 2 处，故有内部分割线。
\begin{tikzpicture}[scale=0.82, semithick,
  dim/.style={<->, >=Stealth, thin},
  lab/.style={font=\small}]
  % ---------- 正视图 ----------
  \begin{scope}
    \draw (0,0) -- (5,0) -- (2,4) -- cycle;
    \draw[dim] (0,-0.35) -- (2,-0.35);
    \node[below] at (1,-0.35) {2};
    \draw[dim] (2,-0.35) -- (5,-0.35);
    \node[below] at (3.5,-0.35) {3};
    \draw[dim] (-0.35,0) -- (-0.35,4);
    \node[left] at (-0.35,2) {4};
    \node[below] at (2.5,-1.0) {正（主）视图};
  \end{scope}
  % ---------- 侧视图 ----------
  \begin{scope}[shift={(7.0,0)}]
    \draw (0,0) -- (4,0) -- (0,4) -- cycle;
    \draw[dim] (0,-0.35) -- (4,-0.35);
    \node[below] at (2,-0.35) {4};
    \node[below] at (2,-1.0) {侧（左）视图};
  \end{scope}
  % ---------- 俯视图 ----------
  \begin{scope}[shift={(0.6,-6.2)}]
    \draw (0,4) -- (5,4) -- (5,0) -- cycle;
    \draw (2,4) -- (5,0);
    \node[below] at (2.5,-0.55) {俯视图};
  \end{scope}
\end{tikzpicture}
